\documentclass[a0,portrait,25pt]{a0poster}
\usepackage{multicol}
\columnsep=100pt
\columnseprule=5pt
\usepackage[svgnames]{xcolor}
\usepackage{times}
\usepackage{graphicx}
\graphicspath{{figures/}{./}} 
\usepackage{booktabs}
\usepackage[font=small,labelfont=bf]{caption}
\usepackage{amsfonts, amsmath, amsthm, amssymb}
\usepackage{wrapfig}
\usepackage{titlesec}
\usepackage{float}

% Custom section formatting mimicking standard poster styles
\titleformat{\section}{\Huge\color{NavyBlue}\bfseries}{\thesection}{1em}{}[{\color{NavyBlue}\titlerule[3pt]}]

\begin{document}
\begin{minipage}[b]{0.20\linewidth}
% Note: Ensure 'fisat.eps' is present in your working directory.
\includegraphics[width=12cm]{fisat.eps}
\end{minipage}
\begin{minipage}[b]{0.78\linewidth}
\begin{center}
\veryHuge \color{DarkRed} \textbf{Meta-WebGuard: A Stacked Ensemble Framework for Real-time Web Threat Detection} \color{Black}\\ [0.5cm]
\Large Rahul K J, Sethulakshmi P A, Sidharth C J, Susan Joy\\ 
\Large Guided by: Dr. Paul P Mathai\\
\Large Department of Computer Science and Engineering, FISAT\\
\end{center}
\end{minipage}
\vspace{2cm}

\begin{multicols}{3}

\section*{Abstract}
Meta-WebGuard is a comprehensive web security extension engineered to identify, intercept, and mitigate cyber threats in real-time. Addressing the failure of traditional static filters against mutated payloads and sophisticated phishing attempts, WebGuard utilizes a multi-layered hybrid framework. By combining $O(1)$ whitelist checking, regex-based static heuristics, and an advanced AI ensemble pipeline consisting of Calibrated SVM, CNN-BiLSTM, and an XGBoost Meta-Learner, the system predicts risk probability with sub-5ms latency. Tested on the Titan-98 Global Mix dataset of over 100,000 URLs, the system achieved a record precision of 99.23\% in detecting zero-day threats, intelligent brand-spoofing, and embedded web vulnerabilities before the user is compromised.

\section*{Problem Statement}
To develop a predictive, low-latency browser guardrail utilizing machine learning to identify phishing, malware, and injection payloads (XSS, SQLi, Path Traversal). Traditional static systems are easily bypassed by entropy-rich payloads, and single-paradigm ML models often suffer from high false-positive rates or overfit to specific domains. The aim is to build a hybrid pipeline integrating statistical, textual, and neural models via Meta-Learning (stacking) to arbitrate between foundation models, providing high-precision, cross-domain threat detection directly in the browser.

\section*{Methodology \& Implementation}

The architectural design of Meta-WebGuard ensures that inference latency remains highly optimized (under 4.5ms), providing seamless, invisible protection via a Flask-based backend server interacting with a React/Vite Chrome Extension.

\textbf{1. The Shield Structure \& Static Heuristics (Layer 1 \& 2):}
Initial layers perform rapid $O(1)$ whitelist checking against known safe internal domains and global high-trust entities. Simultaneously, regex-based heuristic pattern matching instantly flags obvious XSS, SQLi, and Path Traversal attempts without invoking the AI ensemble, saving computational overhead.

\textbf{2. Feature Extraction Pipeline:}
For URLs passing the initial heuristic filters, the system extracts 13 hand-crafted structural features (e.g., Shannon entropy, digit density, URL length, brand spoof indicators). High-dimensional TF-IDF vectors are extracted for lexical NLP representation, and the URL is tokenized and padded into a 550-length sequence for deep learning.

\textbf{3. AI Foundation Models \& Meta-Learner (Layer 3):}
\begin{itemize}
    \item \textbf{Calibrated Support Vector Machine (SVM):} Processes the TF-IDF text representation to detect known lexical signatures, applying `CalibratedClassifierCV` to return probabilities instead of scores.
    \item \textbf{Hybrid CNN-BiLSTM:} A deep sequential model that runs a 1D Convolutional Neural Network (Conv1D) to detect local spatial anomalies, followed by Bi-Directional LSTM layers to capture long-range contextual intent across the URL characters.
    \item \textbf{XGBoost Stacking Arbitrator:} The SVM and neural model probabilities are concatenated with the 13 hand-crafted features. This fused array is passed into an XGBoost tree with over 300 estimators. Utilizing cost-sensitive sample weighting (2.5x penalty for malicious entries), the Meta-Learner yields a robust final threat prediction.
\end{itemize}



\section*{Dataset \& Engineering (Titan-98)}
The ensemble is trained on the \textbf{Titan-98 Global Mix} dataset—a massive fusion of over 100,000 heterogeneous URLs ensuring generalizability.
\begin{itemize}
    \item \textbf{Benign Traffic:} Majestic Million domains and modern benign datasets (to teach the model safe, everyday traffic).
    \item \textbf{Phishing \& Malware:} Active PhishTank campaigns and URLhaus malware drops, introducing zero-day evasive techniques.
    \item \textbf{Injection Logic:} CAPEC-aligned injection datasets containing complex cross-site scripting (XSS) and SQL injection payloads.
\end{itemize}
Crucially, the dataset heavily stratified classes and utilized a 2.5x importance weight on malicious samples to explicitly minimize the False Negative Rate (FNR).

\section*{Experimental Results}
Extensive benchmarking on a 100\% external, unseen URL subset validated the robustness of the stacking architecture against single-layer models. The Meta-Learner reduced False Positives drastically while retaining strong recall on injected payloads.

% Suggestion: Include a confusion matrix here to show true positives vs false positives
\vspace{0.5cm}
\noindent\begin{minipage}{\columnwidth}
\centering
\includegraphics[width=0.85\linewidth]{v6_confusion_matrix.png}
\captionof{figure}{\color{DarkRed} Confusion Matrix demonstrating low False Positives}
\end{minipage}
\vspace{0.5cm}

\vspace{0.5cm}
\noindent\begin{minipage}{\columnwidth}
\centering
\captionof{table}{\color{DarkRed} Master Verification Benchmarks (Mixed Unseen Dataset)}
\renewcommand{\arraystretch}{1.2}
\begin{tabular}{l c}
\toprule
\textbf{Metric Type} & \textbf{Achieved Value} \\
\midrule
Overall Accuracy & 97.70\% \\
Precision (Malicious) & \textbf{99.23\%} \\
Recall (Malicious) & 96.15\% \\
F1-Score & 97.66\% \\
Inference Latency & $<$ 4.5ms \\
\bottomrule
\end{tabular}
\end{minipage}
\vspace{0.5cm}

\vspace{1cm}
% Suggestion: Include the Precision Matrix to prove 99.23% precision assertion
\vspace{0.5cm}
\noindent\begin{minipage}{\columnwidth}
\centering
\includegraphics[width=0.95\linewidth]{v6_precision_matrix.png}
\captionof{figure}{\color{DarkRed} Precision analysis highlighting superior threat detection}
\end{minipage}
\vspace{1cm}

\section*{Conclusions}
\begin{itemize}
    \item Replacing fragile static filters with a multi-layered Hybrid framework protects end-users from sophisticated zero-day campaigns.
    \item Combining TF-IDF semantic scanning with a Deep CNN-BiLSTM sequential analyzer bridges the gap between shallow heuristic rules and powerful contextual ML prediction.
    \item By using XGBoost to cross-validate multiple model outputs, Meta-WebGuard achieves an exceptional $99.23\%$ precision and sub-5ms operational latency.
    \item Cost-sensitive meta-learning successfully proved that complex predictive models can be efficiently operationalized for browser-side cybersecurity guardrails.
\end{itemize}

\section*{References}
\small
\begin{itemize}
    \item D. Demirel et al., "Zero-Shot CNN for Web Attacks," IEEE, 2023.
    \item S. Bensaid et al., "CNN-BiLSTM IDS for SDN Networks," IEEE Transactions, 2023.
    \item F. Zhou et al., "Hybrid Char-SVM \& CNN for Malicious URL Detection," Springer, 2024.
\end{itemize}

\end{multicols}

\vfill
\vspace{0cm}

\begin{center}
\includegraphics[width=0.85\textwidth]{sys_final.png}
\vspace{1cm}
\captionof{figure}{\LARGE \color{DarkRed} Meta-WebGuard System Architecture Workflow}
\end{center}

\end{document}
